\section{Use of generative AI}

Tool: Claude (Anthropic) via Claude Code, used throughout as a coding, analysis and writing
assistant; it wrote and executed code in the student's environment and drafted report text.
In line with the course AI policy, each row gives the prompt (paraphrased in formal language,
preserving its substance), what the tool produced, and what the student changed or decided.
Every number in this report is produced by, and reproducible from, the code in the repository.

{\footnotesize
\begin{longtable}{p{0.025\textwidth}p{0.11\textwidth}p{0.26\textwidth}p{0.26\textwidth}p{0.21\textwidth}}
\toprule
\# & Stage & Prompt (paraphrased) & Output & Student edits and decisions \\ \midrule
\endhead
1 & Planning & Read the assignment handout and its front-page instructions. Without writing code or installing anything, assess whether both questions can be completed, and propose an approach that keeps accuracy while limiting cost, delegating sub-tasks to lower-cost models. & A feasibility analysis for Question 1 (written design) and Question 2 (experiment pipeline checked against the handout's built-in checks). A diagnosis of why the local GPU driver failed to load. A staged plan with a model allocation per sub-task, and a list of decisions for the student. & Reported the deadline extension to Oct 29. Fixed the GPU driver personally by enrolling the Secure Boot key. Chose to manage all GitHub commits and pushes personally. Asked for the Q2.5 and Q2.11 predictions to be drafted by the tool and reviewed by the student before they are recorded. \\ \midrule
2 & Environment & Set up the environment before any experiment, following the approach used in Assignment 1. & A Python 3.12 virtual environment reusing the existing CUDA build of PyTorch 2.9.0, with the matching torchvision 0.24.0, OpenCV, matplotlib and imageio. The Moving MNIST download. A repository skeleton (report folder, this log, the pre-run predictions file). & None. \\ \midrule
\bottomrule
\end{longtable}
}
